% 脚注・割注の横組統合検証。
\documentclass[pdflatex]{jlreq}
\begin{document}
本文中の脚注\footnote{これは脚注である。二行にわたる長い脚注で，
禁則とペアグルーが脚注内でも機能するかを見る。「括弧」や句読点「，。」
も含める。}と割注\warichu{自動位置計算の割注。二行に折り返される長い内容をここに置く。}を含む段落。
割注は割注開き括弧「（」の幅をJFMから測定するため，
版面計算が正しく行われればエラーなく組める。

さらに脚注\footnote{もう一つの脚注。混植 English も数式 $e^{i\pi}+1=0$ を含む。}
を重ねて，脚注番号の採番も確認する。
\end{document}
